\section{Related Work}
\label{sec:related}

\paragraph{Distributed and geo-distributed LLM routing.}
Distributed LLM serving systems increasingly incorporate runtime load and
cache locality into request routing.
Preble jointly considers reusable KV state and computation load when
scheduling prompts across distributed replicas~\cite{preble}, while DualMap
balances cache affinity against load and performs hotspot-aware
rebalancing~\cite{dualmap}.
Recent systems explicitly target geographically distributed serving.
SkyWalker handles cross-region traffic while preserving KV-cache locality and
load balance~\cite{skylb};
GORGO incorporates WAN latency, queueing state, and Prefill/cache cost into
cross-region routing decisions~\cite{gorgo};
and recent work such as Solyx combines serving telemetry with WAN
measurements for per-request cross-site placement~\cite{solyx}.
Although these systems may redirect or rebalance traffic after an initial
choice, RAVEL differs in the control abstraction it exposes:
it explicitly couples \emph{immediate prospective demand accounting} with
\emph{pre-handoff placement revisability}.
An arrived request is accounted at a tentative replica immediately, while
that replica may still change until the request is handed off to the serving
engine.

\paragraph{SLO-aware LLM serving and application scheduling.}
Prior systems also adapt LLM execution to request-level latency and SLO
requirements.
Sarathi-Serve uses Chunked Prefill to reduce Prefill--Decode interference and
control TTFT and inter-token latency~\cite{sarathi}, while DistServe
disaggregates Prefill and Decode and provisions the two phases according to
TTFT and TPOT requirements~\cite{distserve}.
JITServe targets heterogeneous latency-, deadline-, and compound-request
SLOs, scheduling serving capacity from imprecise request information and
refining those estimates as generation progresses~\cite{jitserve}.
AdaServe instead customizes speculative decoding to per-request SLO
requirements~\cite{adaserve}.
Other systems exploit structure across related LLM calls:
Parrot exposes application-level semantic dependencies to coordinate
multi-request execution~\cite{parrot}, and Agentix treats agentic programs as
first-class scheduling entities to optimize program-level end-to-end
latency~\cite{agentix}.
Libra further partitions requests into cooperating micro-requests and combines
global split-point decisions with local SLO-aware batching to balance dynamic
serving load across instances~\cite{libra}.
RAVEL applies the same broad principle of request-specific SLO awareness at a
different control layer: it uses WAN delay and replica serving state to
determine which geographic capacity remains viable for each request before
engine handoff.
Its local Prefill control complements this global decision by bounding
Prefill interference in accordance with latency-sensitive execution
constraints.

\paragraph{Dynamic reassignment and migration.}
Serving systems can also revise placement after an initial assignment.
Llumnix dynamically reschedules requests across model instances and uses live
migration to improve load balancing and isolation, mitigate resource
fragmentation, and differentiate request priorities and SLOs~\cite{llumnix}.
Such migration preserves placement flexibility even after execution state has
materialized.
RAVEL chooses an earlier control boundary: a request remains reassignable
while it is Router-owned, but its placement becomes final after successful
engine handoff.
Thus, RAVEL performs reassignment before changing placement would require
coordination with engine-managed request state, whereas migration systems
retain flexibility after that state already exists.
The two approaches are complementary: pre-handoff replanning can avoid some
state movement, while runtime migration can correct imbalance that emerges
after execution begins.
